\BeleskeID{02-s028}
% element: 02-s028:e01
% element: 02-s028:e02
% element: 02-s028:e03
% element: 02-s028:d01

Četiri prikazana načina dobijanja invertora slede neposredno iz funkcija dvoulaznih NI i NILI kola:
\[
\overline{xx}=\bar x,\qquad
\overline{x\cdot1}=\bar x,\qquad
\overline{x+x}=\bar x,\qquad
\overline{x+0}=\bar x.
\]
Stalna jedinica ne menja I proizvod, a stalna nula ne menja ILI zbir; izlazna negacija u oba slučaja daje invertovanje.

Za pet ulaza kaskada I kola računa najpre \(x_1x_2\), zatim \((x_1x_2)x_3\), pa dodaje \(x_4\) i \(x_5\). Asocijativnost I funkcije omogućava ovu zamenu; analogno važi za ILI. Prva dva ulaza prolaze kroz četiri kola, a poslednji samo kroz završno kolo. Ako sva kola imaju jednako kašnjenje \(t_p\), kritična putanja kaskade ima približno \(4t_p\), dok putanja poslednjeg ulaza ima \(t_p\).
